\documentclass[11pt]{article}
\usepackage[a4paper,margin=25mm]{geometry}
\usepackage[T1]{fontenc}
\usepackage{lmodern}
\usepackage{booktabs}
\usepackage{amsmath}
\usepackage{graphicx}
\usepackage{hyperref}
\usepackage{xcolor}
\usepackage{listings}
\lstset{basicstyle=\ttfamily\small, columns=fullflexible, keepspaces=true,
  breaklines=true, frame=single, framerule=0.3pt, rulecolor=\color{gray}}
\newcommand{\code}[1]{\texttt{#1}}
\newcommand{\us}{\textmu s}

\title{\code{cmnoise.py}: the noise common to the biased channels}
\author{BGOTSCAN}
\date{2026-10-15}

\begin{document}
\maketitle

\section{Purpose}

The charge of the summed waveform over the 5~\us{} window carries a noise
of 66\,000 ADC\,ns at 41~V, while one live channel carries 2400: if the
52 live channels were independent the sum would carry $2400\sqrt{52}
\approx 17\,000$. Most of the noise of the sum is therefore common to the
live channels, and the twelve dead channels, which show pure white noise,
do not see it (\code{chanstat\_doc.tex}, \code{ANALYSIS.md} 2026-10-15).
Two sources are on the table: the trigger logic, which fires about five
times per event at 41~V, recorded or not, and kicks the biased channels
as the excursion at the trigger time shows; or a continuous noise common
to the biased channels, from the bias line or elsewhere. The first gives
discrete kicks in the tail of the events at the trigger rate; the second
a continuous, Gaussian noise. This script isolates the common component
in the tail of the events, shows it, measures how common it is and among
which channels, takes its spectrum, and searches the tail for kicks with
the shape of the one at the trigger time. It reads the raw HDF5 file with
h5py, like \code{chanstat.py}, and the per-event classes of
\code{offline.py} to keep only clean events.

\section{Usage}

\begin{lstlisting}
python3 src/py/cmnoise.py RUN [--events 0-299] [--tail 12] [--smooth 125] [--dead LIST]
        [--offline FILE] [--kick 16] [--nsigma 5] [--traces 4] [--save NAME]
\end{lstlisting}

\begin{center}
\begin{tabular}{l l p{80mm}}
\toprule
option & default & meaning \\
\midrule
\code{RUN} & & run name or path of the raw file \\
\code{--events a-b} & \code{0-299} & events read (0-based) \\
\code{--tail T} & 12 & the tail: the last $T$~\us{} of the record \\
\code{--smooth N} & 125 & running-mean width in samples (1~\us) of the slow component \\
\code{--dead LIST} & the twelve of \code{chanstat\_doc.tex} & dead channels \\
\code{--offline FILE} & \code{pass/RUN\_offline.h5} & per-event classes; without it every event is used \\
\code{--kick K} & 16 & half-width of the kick template in samples \\
\code{--nsigma S} & 5 & threshold of the kick search, in units of the filtered noise \\
\code{--traces N} & 4 & events drawn in the trace figure \\
\code{--save NAME} & & save the figures as \code{src/py/plots/NAME\_cmnoise\_*.png} \\
\bottomrule
\end{tabular}
\end{center}

\section{Definitions}

A record is one event as the digitizer stores it: 3000 samples at 8~ns of
the 64 channels, starting 1.6~\us{} before the trigger; ``event'' counts
them, ``record'' refers to the waveform window. Every waveform is taken minus the WebDump baseline of its channel (the
attribute \code{baselines}, one value per channel for the run), so that
the offset of a record with respect to the run is kept. The events used
are the REAL events without pile-up of \code{offline.py} whose maximum
lies more than 4~\us{} before the tail, so that the scintillation pulse
is below 3\% of its height throughout the tail.

\paragraph{The two means.} Per event, over the tail, the mean of the live
channels $L_i$ and the mean of the dead channels $D_i$, sample by sample,
and their running means over $N$ samples $\bar L_i$, $\bar D_i$. For white
noise of $\sigma$ counts per channel, the rms of a mean of $n$ channels is
$\sigma/\sqrt n$, and of its running mean $\sigma/\sqrt{nN}$: with
$\sigma = 4$, 0.55 and 0.05 counts for the 52 live channels, 1.15 and
0.10 for the 12 dead ones. What $L$ has above these numbers is common to
the live channels; $D$ is the control that it is not in the digitizer.

\paragraph{Traces.} Figure \code{traces}: for the first \code{--traces}
events, $L_i$ and $D_i$ as points and $\bar L_i$, $\bar D_i$ as lines
against time, one panel per event, the same vertical scale.

\paragraph{How much.} Figure \code{rms}: per event, the rms of $L$, $D$,
$\bar L$ and $\bar D$ over the tail, each minus its own mean over the
tail, as histograms over the events with the white-noise values marked;
and the pooled distribution of the $\bar L$ samples with a Gaussian of
the same rms, for the shape. The terminal prints the medians, the
correlation between $\bar L$ and $\bar D$, and the scatter of the mean of
$L$ over the tail between events (the event-to-event baseline scatter,
47 counts on the sum at 41~V in \code{ANALYSIS.md}).

\paragraph{Among whom.} The correlation matrix of the 64 channels: for
each event the 64 running means over the tail, each minus its mean, and
the covariance accumulated over the events, normalised to a correlation.
Figure \code{corr}: the matrix, channels in order, dead channels marked.
The terminal prints the mean correlation live--live, dead--dead,
live--dead, the mean correlation within and across the four $4\times4$
SiPM matrices of the map (a block structure points at the hardware that
shares the noise), and the fraction of the live-channel variance in the
first principal component with the spread of its loadings (a uniform
loading is a noise common to all with the same weight).

\paragraph{Spectrum.} Figure \code{spectra}: the power spectra of $L$,
$D$ and $L - D$ over the tail, averaged over the events, against
frequency.

\paragraph{Between events.} The offset of each record from the run
baseline, the mean over the tail per channel: its scatter between events
for the live mean, the dead mean and a single channel, the correlation
of the live and dead means, and the mean correlation between channels
(live--live, dead--dead, live--dead). A constant offset over the record
cancels in the charge, which takes its baseline from the same record;
this is the quantity that was read as a baseline scatter following the
trigger rate (\code{ANALYSIS.md}, 2026-10-15 morning).

\paragraph{Kicks.} The median over the live channels of the record
(robust against the lit channels), averaged over the events around the
trigger sample, and the dead-channel mean averaged the same way, show
what the biased and the unbiased channels do when the trigger fires. A
first version took the extremum of that average as a kick template; it
is not a kick but a step, the small pulse that the median channel
carries, so the search uses a boxcar: the running mean of the live
median of the tail over $2K + 1$ samples (\code{--kick} $K$, 264~ns by
default), and a kick is a local extremum above \code{--nsigma} times its
clipped rms with no larger one within a boxcar length. Figure
\code{kick}: the two averages around the trigger, the per-event
threshold, the running mean of one event with the threshold, and the
amplitudes found. The terminal prints the step at the trigger, the
threshold in counts per channel and on the sum, and the kicks per event
and per \us{} against the expectation from the board trigger rate (read
from \code{pass/RUN\_offline.json} when present) times the tail length:
2.7 per 12~\us{} tail at 41~V, 0.03 at 39~V.

\section{Findings (2026-10-15, 195K\_41V\_th20 and 197K\_39V\_th20, 300 events each)}

\begin{itemize}
\item \emph{No kicks.} Any excursion of the live median above 0.94
  counts per channel (49 on the sum) lasting 264~ns would be found; 0 in
  236 clean events at 41~V against 637 expected if the trigger logic
  kicked the channels at each of its 225~kHz firings. The trigger
  hypothesis is excluded at that size. The feature at the trigger time
  is a step of $+3.3$ counts per channel in the live median, $+0.05$ in
  the dead mean: the pulse of the median channel, light and not pickup;
  the ``excursion of 200 counts on the sum at the trigger'' of
  \code{selection\_doc.tex} is this step times 52.
\item \emph{Within the record} the common noise among the live channels
  is weak: pairwise correlation 0.03 of the 1~\us{} running means, with
  no block structure (the four SiPM matrices 0.02 to 0.05, within and
  across alike) and a first principal component holding 5\% of the
  variance. The dead channels show none (0.00). Weak per pair, it still
  dominates the sum: $1 + 51 \times 0.03 = 2.5$, so about 60\% of the
  slow variance of the sum is this common part and 40\% the independent
  per-channel slow noise.
\item \emph{The per-channel slow noise is in every channel}, dead or
  live: the 1~\us{} running mean of the live mean has an rms of 0.117
  against 0.047 for white noise, the dead mean 0.181 against 0.098, the
  same excess once the number of channels is accounted for, and the
  spectra of both rise below 0.3~MHz the same way. It is in the
  front-end or the digitizer, not in the bias or the SiPMs.
\item \emph{Nothing depends on the bias or the trigger rate.} Every
  number above is the same at 41~V (225~kHz) and at 39~V (2.7~kHz):
  running-mean rms 0.117 and 0.118, correlations 0.03 and 0.03,
  spectra identical. The baseline scatter that fell as the square root
  of the trigger rate (47 counts at 41~V to 5.6 at 39~V) compared a pass
  with the dead-channel common mode to passes with the interpolated
  median: with the same method the pass files give 47 and 29 (clipped
  32 and 28), and this script 44 and 48 on clean events.
\item \emph{Between events} every channel's offset scatters by 1.3
  counts, dead or live, with a correlation of 0.4 to 0.5 between any two
  channels: a board-wide offset of about 0.9 counts common to the 64
  channels plus a per-channel part of about 1 count. It cancels in the
  charge.
\item \emph{An 18~MHz line} in the live mean, absent from the dead
  mean, at both biases: a coherent pickup in the biased channels, far
  above the frequencies the charge integrates over.
\end{itemize}

Consequence: the noise of the charge is the slow noise within the
record, two parts, a per-channel one from the front-end and a weakly
common one among the biased channels, neither tied to the trigger. The
software handle is per channel (a baseline model per channel over the
window, as the per-channel noise is 2400 ADC\,ns against 700 for white
noise) and fewer channels in the sum; the hardware question is the
front-end's low-frequency noise.

\section{Examples}

\begin{lstlisting}
python3 src/py/cmnoise.py 195K_41V_th20 --save cm195   # docs/figs/cmnoise_195K_{traces,rms,corr,spectra,kick}.png
python3 src/py/cmnoise.py 197K_39V_th20 --save cm197
\end{lstlisting}

\end{document}
